\subsection*{MT-M29. The determinant of the graded arithmetic spectral module}
Fix one unitary Hecke character and its pole-removed entire
completion \(\Xi\). Work first in the centered, actual wide-strip
Hilbert quotient of MT-M24/MT-M26, with \(a=1/2\) and
\([0,1]\subset1/2+\operatorname{int}I\). Its definition by adelic
summation and Hilbert closure is unconditional. Write \(U_y\)
for norm translation by \(\tau(e^y)\), and \(D\) for its
infinitesimal generator. Its dense smooth quotient is the
centered version of \(Q_\eta\).

There is a canonical algebraic submodule \(V_{\mathrm{fin}}\):
the vectors contained in finite-dimensional \(U_y\)-invariant
subspaces. It is dense in the smooth quotient and in the
specified Hilbert quotient. Its primary decomposition and
canonical filtration are
\[
 V_{\mathrm{fin}}=\bigoplus_\rho V_\rho,\qquad
 V_\rho\simeq\mathbf C[u]/(u^{m_\rho}),\qquad
 D|_{V_\rho}=(\rho-\tfrac12)+u,\qquad
 F^jV_\rho=(D-\rho+\tfrac12)^jV_\rho .
                                                               \tag{M29.1}
\]
The isomorphism uses the local Mellin coordinate
\(u=z-(\rho-1/2)\). Form the associated graded
\[
       V_{\mathrm{gr}}
          =\bigoplus_\rho\bigoplus_{j=0}^{m_\rho-1}
                         F^jV_\rho/F^{j+1}V_\rho.             \tag{M29.2}
\]
This construction is canonical and functorial at each finite
spectral stage for maps intertwining norm translation. It
retains \(m_\rho\) dimensions at \(\rho\); it replaces the
nilpotent extensions by scalar actions.

Assume now GRH for this function. It is supplied by \(A_D\)
for a Dirichlet function over \(\mathbf Q\), or by the
conditional input \(M\) for the general Hecke character.
Give the one-dimensional graded lines in (M29.2) the metric
in which the class of \(u^j\) has norm one, and make them
orthogonal. The resulting Hilbert completion
\(\mathcal H_{\mathrm{gr}}\) carries the unitary norm action
\[
        U_y^{\mathrm{gr}}=e^{iyT},\qquad
        T|_{(\rho,j)}=\gamma,\qquad \rho=\tfrac12+i\gamma.
                                                               \tag{M29.3}
\]
Every integrated \(\mathcal E\)-test operator is trace class,
with the same trace as on the actual wide-strip quotient:
\[
 \operatorname{Tr}_{\mathcal H_{\mathrm{gr}}}C_f^{\mathrm{gr}}
       =\sum_\rho m_\rho\,\mathcal Mf(\rho-\tfrac12)
       =\operatorname{Tr}_{Z_{1/2,I,s}} C_f .                  \tag{M29.4}
\]
The grading and this new Hilbert completion change the
representation when a multiple zero occurs; they do not
equip the original Jordan blocks with a positive invariant
metric.

Suppose in addition that
\[
 \Xi(\tfrac12+z)=\epsilon\Xi(\tfrac12-z),\qquad
       \epsilon\in\{1,-1\}.
                                                               \tag{M29.5}
\]
Delete the central spectral block, then take the
positive-frequency spectral subspace
\(\mathcal H_+=1_{(0,\infty)}(T)\mathcal H_{\mathrm{gr}}\).
On it,
\[
                          B=T_+^{-2}                         \tag{M29.6}
\]
is a bounded positive trace-class operator. If
\(r=\operatorname{ord}_{s=1/2}\Xi(s)\), set
\[
 H(t)=t^{-r}\Xi(\tfrac12+it),\qquad
              \mathfrak D(w)=H(\sqrt w)/H(0).
 \]
Then the operator from the actual arithmetic module has
\[
                       \det(1-wB)=\mathfrak D(w).             \tag{M29.7}
\]
This is the AG-DET operator, now identified as a functional
calculus of the norm generator on the graded finite
spectral module. Its numerical multiplicities are precisely
the original zero orders.

\begin{proof}
\emph{The finite spectral core comes from the actual quotient.}
MT-M18 and MT-M24 identify the smooth quotient with its full
zero jets. For a finite set \(S\) of zeros, there is a unique
class with any prescribed jets at \(S\) and zero jets at
every other zero: its existence follows from the finitely
many block lifts of MT-M18, and uniqueness from the exact
closed-ideal equality. These classes form a finite-dimensional
subalgebra \(V_S\), canonically isomorphic to
\[
                     \prod_{\rho\in S}
                              \mathbf C[u]/(u^{m_\rho}).
                                                               \tag{M29.8}
\]
It has a canonical unit \(e_S\), whose local class is one
at \(S\) and zero elsewhere. Multiplication by \(e_S\) is
the finite-jet quotient projector. Thus \(V_S\) is both
a subspace and a finite quotient of the actual smooth
arithmetic quotient; no choice of interpolation
representatives survives in (M29.8).

These finite projectors extend boundedly to the actual
wide-strip Hilbert quotient: its finite jet restrictions
are continuous by MT-M26, and the finitely many lifted
classes have finite Hilbert norm. This proves that \(V_S\)
is a genuine finite spectral submodule there as well.
The finite unions of the MT-M18 blocks exhaust all zeros,
and their partial sums converge in the smooth quotient.
Hence \(\bigcup_S V_S\) is dense in that quotient and,
by its dense inclusion, in the Hilbert quotient.

Norm translation acts in the Mellin transform by
multiplication by \(e^{zy}\) in the centered variable.
Consequently its action on the \(\rho\)-factor in (M29.8)
is multiplication by
\[
                  e^{(\rho-1/2+u)y}
                   =e^{(\rho-1/2)y}e^{uy}.
 \]
Differentiating this finite matrix exponential gives
the generator in (M29.1).

It remains to verify that this core is intrinsic, rather
than merely a specified collection of finite subspaces.
Let \(W\) be any finite-dimensional invariant subspace of
the Hilbert quotient. The continuous one-parameter group
on \(W\) is a matrix exponential; write \(D_W\) for its
matrix generator, and let \(P\) be its minimal polynomial.
Every finite block restriction intertwines the norm
group, and therefore intertwines its generator on \(W\).
On a local factor with \(P(\rho-1/2)\ne0\),
\[
                 P((\rho-1/2)+u)
 \]
is invertible, because its constant term is nonzero.
Thus the restriction of every vector of \(W\) to this
factor is zero. Only the finitely many zeros whose
centered parameters are roots of \(P\) can remain.
If \(S\) is that finite set, subtracting the finite
projector \(e_S\) kills all jets of a vector in \(W\).
The exact Hilbert kernel statement of MT-M26 forces
that difference to be zero. Therefore \(W\subset V_S\).
Conversely every \(V_S\) is finite dimensional and
invariant, so \(V_{\mathrm{fin}}=\bigcup_S V_S\).
This establishes the intrinsic description.

\emph{The filtration and its functoriality.}
For the primary summand \(V_\rho\), put
\(N_\rho=D-(\rho-1/2)\). Its image powers
\(N_\rho^j V_\rho\) are the filtration in (M29.1);
in the coordinate algebra they are \((u^j)\).
Any finite-stage intertwining map commutes with \(D\).
It preserves the primary decomposition and sends
\(N_\rho^j V_\rho\) into the corresponding image power
in its target. It therefore induces a map on every
graded quotient, compatibly with compositions.
The inclusions as \(S\) grows have the same property.
This proves the claimed algebraic functoriality.
No boundedness assertion for arbitrary infinite-stage
intertwiners is required.

On \(F^jV_\rho/F^{j+1}V_\rho\), the nilpotent part
\(N_\rho\) acts as zero, so \(U_y\) acts by
\(e^{(\rho-1/2)y}\). Each of these quotients has
dimension one, and there are exactly \(m_\rho\)
of them. The distinguished vector \([u^j]\)
uses the local unit in (M29.8) and the specified
global Mellin coordinate; it is independent of
any chosen theta-function representative.
Declaring these vectors orthonormal defines
the metric stated above.

\emph{The Hilbert action and its trace.}
Under GRH the centered parameter is \(i\gamma\).
Thus every norm-translation scalar on the associated
graded has modulus one. The Hilbert completion is
explicitly
\[
             \ell^2\{(\rho,j):0\le j<m_\rho\}.
 \]
The norm group acts by the scalars \(e^{i\gamma y}\).
It is strongly continuous, first on finite-support
vectors and then on all vectors by their density
and the unit norm of the action.
Its self-adjoint frequency generator is multiplication
by \(\gamma\), with domain
\[
       \left\{(a_{\rho,j}):
                    \sum_{\rho,j}\gamma^2|a_{\rho,j}|^2
                                      <\infty\right\}.
 \]
Self-adjointness follows directly: the adjoint of
this multiplication operator has exactly that
domain, since testing against the finite-support
coordinate vectors determines its coordinates.
This proves (M29.3).

For \(f\in\mathcal E\), the integrated action is
diagonal with entry \(\mathcal Mf(i\gamma)\)
on each of the \(m_\rho\) graded lines. The transform
decreases faster than every power of \(|\gamma|\).
The fixed-function local zero count gives
\(\sum_\rho m_\rho(2+|\gamma|)^{-2}<\infty\).
Hence this diagonal operator is trace class, with
trace and trace norm obtained by summing its
entries and their absolute values, respectively.
MT-M26 gives the identical trace on the actual
wide-strip Hilbert quotient. This proves (M29.4).
The proof also gives continuity into trace norm
for the \(\mathcal E\) test topology.

The central subspace has dimension \(r<\infty\).
After it is removed, positive zero ordinates
are bounded away from zero, since zeros of a
nonzero entire function have no finite accumulation
point. Thus \(T_+^{-2}\) is bounded and positive,
with eigenvalues \(\gamma^{-2}\), one for each of
the \(m_\rho\) graded lines at each positive
ordinate. Their sum converges by the same
zero count. This proves (M29.6), including trace
class and compactness. If there are no positive
zeros the space is zero and its determinant is one.

By (M29.5), the noncentral ordinates occur in
pairs \(\gamma,-\gamma\) with equal multiplicities.
AG-DET, whose complete factorization proof
applies to this entire order-at-most-one
completion, states precisely
\[
        \mathfrak D(w)=
          \prod_{\gamma>0}(1-w/\gamma^2)^{m_{1/2+i\gamma}}.
 \]
The spectrum of (M29.6), with its just-proved
graded multiplicities, gives this same locally
uniformly convergent Fredholm product.
This proves (M29.7) and the comparison with
AG-DET. It uses no assertion about individual
Hermite vectors forming a Schauder basis.
\end{proof}

\paragraph{The ordinary trace-GNS completion is different.}
The distinction is essential when zeros have multiplicity.
On the centered convolution algebra \(Q\), let
\[
       \tau(f)=\operatorname{Tr}_{Z_{1/2,I,s}}C_f,\qquad
       f^*(x)=\overline{f(-x)}
 \]
in the fixed-character radial notation. The Hecke
functional equation makes the centered ideal invariant
under this involution. Under the same GRH input,
\[
 \langle[f],[g]\rangle_\tau
          =\tau(f*g^*)
          =\sum_{\rho}m_\rho\,
                    \mathcal Mf(i\gamma)
                        \overline{\mathcal Mg(i\gamma)} .
                                                               \tag{M29.9}
\]
Its radical is exactly the ideal of classes whose
values, without their higher derivatives, vanish at
every zero. Indeed all weights in the sum for
\(f=g\) are positive, so zero norm is equivalent
to all those values vanishing. Finite interpolation
from (M29.8) gives every finitely supported list
of values. The completion of the quotient by
this radical is therefore
\[
       \mathcal H_{\mathrm{GNS}}
             \simeq\ell^2(\{\text{distinct }\rho\},m_\rho).
                                                               \tag{M29.10}
\]
The evaluation map is continuous from \(Q\):
rapid Mellin decay and the local zero count
bound the squared norm by finitely many
\(\mathcal E\) seminorms, and it vanishes on
the closed ideal. Thus the density assertion
in (M29.10) is a statement about the actual
arithmetic quotient, not only an abstract
diagonal model.

Multiplication by a test transform on (M29.10)
has one Hilbert dimension per distinct zero.
The normalized coordinate vector at \(\rho\)
is \(m_\rho^{-1/2}e_\rho\); its diagonal matrix
entry is still \(\mathcal Mf(i\gamma)\).
Consequently its ordinary Hilbert trace counts
each distinct zero once. The integer \(m_\rho\)
in (M29.9) is a weight in the spectral measure,
not a Hilbert spectral multiplicity.
It follows that the positive-frequency inverse
square on this GNS space has ordinary
Fredholm determinant
\[
                \prod_{\substack{\gamma>0\\\text{distinct}}}
                                      (1-w/\gamma^2).         \tag{M29.11}
\]
It agrees with (M29.7) if and only if every
positive zero is simple. This follows also
by comparing the orders of the isolated
zeros of the two entire determinants.
The associated graded construction (M29.2)
is what preserves the multiplicities required
by the arithmetic trace and the completed
function.

There are actual number-field examples in which
this distinction matters. For
\(E=\mathbf Q(\zeta_3,\sqrt[3]{2})\), NT-L17
proves unconditionally that \(\zeta_E\) has
infinitely many zeros of multiplicity at
least two. Under \(M\) these lie on the
critical line; infinitely many are
noncentral, and the selfdual symmetry
gives infinitely many positive ordinates.
Thus the ordinary GNS determinant (M29.11)
does not equal the pole-removed completed
\(\zeta_E\) determinant, whereas the
graded determinant (M29.7) does.
This example uses no simplicity conjecture.

\paragraph{The established family and the remaining geometric gate.}
For a selfdual primitive Dirichlet function,
including the separately pole-removed zeta
function, all of the positive conclusions
above use the designated theorem input
\(A_D\), without \(M\). A complex Dirichlet
character and its conjugate are handled
by the finite direct sum of their actual
arithmetic modules and the product of
their completions. Take the associated
graded in each summand, then the positive
frequency subspace of their direct sum.
The product completion is selfdual, its
zero multiplicities are the sums of the
summand multiplicities, and the resulting
Fredholm determinant is its AG-DET
determinant. Each of these assertions
follows from finite direct sums and the
already proved trace/product formulas.
The same construction applies under \(A_D\)
to a specified Hecke factor whose completed
function is a finite product of Dirichlet
completions with the same center, once
that classical factorization is supplied.
For general Hecke factors, and for their
selfdual finite products, the location
input remains conditional on \(M\).

This advances the comparison requested in
AG-Q08 and in the commentary after AG-DET:
the determinant operator is now obtained
from the canonical filtration of the
locally finite part of the actual
adelic summation quotient. Its positivity
still uses the stated zero-location
theorem. It is not an independently
prime-defined positive cohomology or
an identification with the
Gan--Gross--Prasad period determinants
in Alp\"oge's \emph{Remark}, steps~4
and~7 (PDF pp.~3--4). Those period and
integral-structure comparison maps
remain separate.

The precise Masullo source route is
different from an assertion of positivity
of the full Jordan representation.
Section~1.4.6, original PDF p.~30,
explicitly does not invoke Weil
positivity; \S1.4.8--1.4.10,
pp.~31--32, places the main burden
on computing the same geometric
trace intrinsically in arbitrarily
narrow strip completions.
Theorem~12.5.1.1, pp.~519--520,
equations (12.5.1.1)--(12.5.1.4),
asserts the comparison with the
zeros inside the chosen strip for
tests vanishing on \(C_K^1\).
Original pages~30,31,519,520
were rendered and read; the
complete displayed proof of
Theorem~12.5.1.1 was read as
a dependency map, not as an
audit of its cited upstream
kernel theorems. Equation
(10.3.1.17), p.~450, was
also inspected in the original
page image: its triangular
derivative block is consistent
with the local algebra above,
without supplying the missing
uniform orbital comparison.

MT-M28 proves the narrow-strip
trace class under \(M\). It
does not prove that geometric
comparison independently of
zero location. The remaining
gate is the zero-independent
geometric trace identity, or
an independent positive
cohomological/period comparison.
The new graded determinant
comparison closes the
particular module-identification
step and retains the full
numeric multiplicity, while
making the loss of nilpotent
extension data explicit.

\paragraph{Run locators.}
This new comparison is proved in
\texttt{GRH\_NIGHT/complexity},
\texttt{work/root/graded\_adelic\_determinant.tex}.
It connects MT-M18/MT-M24/MT-M26
to AG-DET and to the inherited
\texttt{ALPOGE\_BRIDGES} rows
BR-TR-01/02, the SA23
\texttt{geo-H-lidskii}/\texttt{geo-H-genus}
rows, and cote~1 scans~79--80
and~111--112 as recorded in
the supplied SA23 intake.
Those archive locators retain
their inherited status.
